\section{The range \texorpdfstring{$D\geq875$}{D >= 875}}\label{app:global}
We now apply the radial estimate to the part of the proof preceding the
finite comparisons of Section~\ref{sec:completion}. The argument has two
parts. On an unbounded interval the arithmetic estimate grows faster
than the size estimate, with an explicit positive difference at the left
endpoint. Below that endpoint we keep the integer parameters fixed on
each of finitely many adjoining intervals.

\subsection{Parameters and geometry on the unbounded interval}
Set
\[
 D_0=5\,760\,000,\quad X=(D-3)/4,\quad X_0=(D_0-3)/4,
 \quad K=31,\quad c_R=23/10,\quad e=29/500.
\]
For $D\geq D_0$ choose
\begin{gather*}
 S=\lfloor c_RX/K\rfloor,\quad R=KS,\quad J=\lfloor eR\rfloor,\\
 T=106\,838,\quad R_0=KT=3\,311\,978,\quad J_0=192\,094.
\end{gather*}
Thus $S\geq T$, $R\geq R_0$ and $J\geq J_0$.
Write $J_*=J_0+1$, $M_0=R_0(T+1)/2$, and put
\begin{align*}
 c_-&=c_R-K/X_0,& e_-&=e-1/R_0,\\
 r_0&=\frac{4\sqrt{1.02K/e_-}}{\sqrt{2\pi T\sqrt{D_0}}}
            +\frac{1+\sqrt{1+4/D_0}}T,&
 \gamma_0&=1+r_0+800e/T.
\end{align*}
These give $\gamma_0<1.002779879<1.02$.
Choose $n$ by rounding $(1+r_0)RS/J$ strictly upward to a multiple of
$800$. Then
\begin{equation}\label{eq:global-node-bounds}
 \frac{RS}{J}<n\leq\gamma_0\frac{RS}{J}.
\end{equation}
The disk bound gives $M_y\leq4\rho/\sqrt D+1$, and substitution of
$n\leq1.02RS/J$ shows $M_y/S\leq r_0$. The selected $n$ therefore
satisfies $nJ>R(S+M_y)$. Indeed, the definition of $n$ implies $n\leq\gamma_0RS/J<1.02RS/J$, which
justifies the radius estimate used to bound $M_y$.

For later use put
\begin{align*}
 \theta_-&=398891129/800000000,&\theta_+&=(1+1/T)/2,\\
 \nu_-&=c_-/(Ke),&\nu_+&=\gamma_0c_R/(Ke_-),\\
 \eta_+&=\frac{\sqrt{D_0+4}/4}
                  {\sqrt{\nu_-X_0\sqrt{D_0}/(2\pi)}},\\
 C_-&=\tfrac12\log(\nu_-/\pi),&
 C_+&=\tfrac12\log\left(\nu_+\sqrt{1+3/(4X_0)}/\pi\right)
                                      +\log(1+\eta_+).
\end{align*}
Equation~\eqref{eq:global-node-bounds} gives
\begin{equation}\label{eq:global-geometry}
 \theta_-\leq\theta\leq\theta_+,\quad
 \nu_-X\leq n\leq\nu_+X,\quad \eta\leq\eta_+,
 \quad \tfrac34\log X+C_-\leq\log\rho\leq\tfrac34\log X+C_+.
\end{equation}
For the first inequality use $\theta_-<1/(2\gamma_0)$.
For the third, the square of the bound for $\eta$ is a constant times
$(D+4)/((D-3)\sqrt D)$, which decreases for $D>4$.
The last inequality follows by writing $D=4X+3$ in the disk formula.
In particular $\eta_+<0.022619737$.
The remaining resultant hypothesis holds as well: $J<R$, and
\[
 \frac{2\rho+1}{R}
 \leq\frac{2e^{C_+}}{c_-}X^{-1/4}+\frac1{c_-X}<1.
\]
The right side decreases and is less than one at $X_0$.
Thus $\max(J,M_x)\leq R+KM_y$, and the resultant criterion supplies
a nonzero maximal minor throughout this interval.

\subsection{A uniform arithmetic lower bound}
Only primes at most $X$ are used in this part of the argument. They are
all inert by Lemma~\ref{lem:small-inert}. We give the lower bound
explicitly, including its dependence on the anchor $T$, so that its
validity as $D$ increases is visible.

For an integer $q\geq4$, put
\[
 E_T(q)=\left[\frac1{Kq}-\frac3{2R_0}+
             \frac3{2R_0q}-\frac{q}{2\sqrt3KT^3}\right]_+.
\]
The floor-sum estimate for $T(S,q)$ gives
$2KT(S,q)/W\geq E_T(q)$ when $S\geq T$.
One can obtain it by expanding the quadratic sum
\eqref{eq:T-floor}; the only signed periodic term is bounded using
$|t(1-t)(1-2t)|\leq1/(6\sqrt3)$, $0\leq t\leq1$.
For $4\leq q\leq S$, the positive term $q/(2R_0^2)$ absorbs the
Bernoulli remainder in the displayed expression; for $q>S$ the sum is zero.
Dropping the remaining nonnegative terms gives the displayed expression.
It increases when $T$ increases, since its terms depending on $T$ are
$-3(1-1/q)/(2KT)-q/(2\sqrt3KT^3)$.

For a prime $p<K$, write $a_p=K-p$ and define
\[
 W_p^{(0)}=\frac{1-T^{-2}}K+
                      \sum_{p^{2l}\leq T}E_T(p^{2l}),
 \qquad C(v)=v^2(3-2v),\quad e_* =J_*/R_0.
\]
For each $v\in\{0,1/200,\ldots,1\}$, put
\[
 v_-=(v-1/T)_+,\qquad q_v=\frac{a_p}{K}v\min(1,v+1/T).
\]
Define the corresponding row gains and column loss by
\begin{align*}
 A_p(v)&=\frac{a_p}{Kp^2}[C(v_-)-3/T]_+
          +\frac{a_p}{Kp^3}[1-C(v)-6/T]_+\\
 &\quad+\min\left(0,\frac{3a_p(K+p-1-2(p^2-1))}{2K^2p^2T}\right)
                  -\frac{3a_pp^2}{2K^2T^2},\\
 B_p(v)&=\left[\frac{a_p}{K^2}
     \left\{\left(\frac1p-\frac1{p^2-1}\right)v_-^3-\frac3{2T}\right\}
                                                        \right]_+,\\
 C_p(v)&=3(1+1/T)q_v(1-q_v\theta_-/2)\frac{p-1}{p}e_*\\
 &\quad+\boldsymbol1_{v>0}
       \frac{3(1+1/T)\lfloor p^2/4\rfloor}
                  {2p\theta_-(J_0/R_0)R_0^2}.
\end{align*}
The resulting small-prime weight is
\[
 W_p=W_p^{(0)}+
       \max\left(0,\max_v\{[A_p(v)+B_p(v)]_+-C_p(v)\}\right).
\]
These are lower bounds for the small-prime row weights in
\eqref{eq:wkeep}. To see the direction of each estimate, take
$L=\lfloor vS\rfloor$ there, so that $v-1/S\leq L/S\leq v$.
The sums of $t(S-t)$ in the two parts of the triangle, after
normalization, have limiting values $C(v)$ and $1-C(v)$; summing
$t$ and $t^2$ bounds their errors by $3/S$ on either part.
The constant term of the strip estimate contributes at least its
negative part. The incomplete-strip loss is at most
$3a_pp^2/(2K^2S^2)$. For the extra odd-power sum, integral comparison
of the floors gives the lower bound in $B_p$.
Finally, the proportion of selected rows is at most $q_v$ and
$q(1-q\theta/2)$ increases for $0\leq q\leq1$; the column estimate
therefore gives $C_p$. Replacing $S$ by $T$ and $\theta$ by $\theta_-$ only weakens these bounds.
For the final incomplete-column remainder, the denominator is bounded using
$JR\geq J_0R_0$, which follows directly from
$J\geq J_0$ and $R\geq R_0$.

Before replacing $S$ by $T$, the corresponding lower gains increase with
$S$, while the losses decrease: $(v-1/S)_+$ increases,
$(a_p/K)v\min(1,v+1/S)$ decreases, and $C$ is increasing on $[0,1]$.
Thus each fixed choice of $v$ gives a uniform lower bound. Taking the
maximum over the grid preserves this bound; any additional tower terms
are nonnegative.

Write $L_0=\log X_0$, $u_-=1/c_R$, $u_+=1/c_-$, and
\[
 E_M=-\gamma_{\rm Euler}-\sum_p\frac{\log p}{p(p-1)}.
\]
The explicit estimates of Rosser and Schoenfeld
\cite[Theorems 4 and 6]{rosserschoenfeld}, followed by partial summation,
give, for $x\geq X_0$,
\begin{align}
 \sum_{p\leq x}\frac{\log p}{p}
       &\geq\log x+E_M-\frac1{2\log x},\notag\\
 \vartheta(x)&\leq x\left(1+\frac1{2\log x}\right),\notag\\
 \sum_{K\leq p\leq x}p\log p
       &\geq\frac{x^2}{2}
         \left(1-\frac3{2\log x}-\frac1{3\log^2x}\right)-160000.
                                                    \label{eq:global-prime-sums}
\end{align}
For the last inequality, integrate the bound for $\vartheta$ from $563$
and use
\[
 \frac{d}{dt}\left(\frac{t^2}{2\log t}
                         +\frac{t^2}{3\log^2t}\right)
 \geq\frac{t}{\log t}\qquad(t\geq563).
\]
The constant $160000$ also covers the omitted initial interval.

The large-prime row weight has the lower bound
\[
 \frac{S_r(R,K,p)}{KW}
 \geq\frac1p-\frac3{2R}+\frac{3K}{2Rp}
            +\frac{p}{2R^2}-\frac{3K}{2R^2}
            -\frac{p^2}{12\sqrt3R^3}.
\]
The two terms containing $3K$ have nonnegative sum when $p\leq X<R$.
Also $p^2\leq Xp$. Applying \eqref{eq:global-prime-sums} therefore
shows that the normalized arithmetic term is at least
\begin{equation}\label{eq:global-arithmetic-growth}
 A_0+\log(X/X_0),
\end{equation}
where
\begin{align}
 A_0={}&\sum_{p<K}W_p\log p
    +\sum_{\substack{p\geq K, l\geq1\\p^{2l}\leq T}}
                                     E_T(p^{2l})\log p\notag\\
 &+L_0+E_M-\frac1{2L_0}-\sum_{p<K}\frac{\log p}{p}
       -\frac3{2c_-}\left(1+\frac1{2L_0}\right)\notag\\
 &+\left(\frac{u_-^2}{2}-\frac{u_+^3}{12\sqrt3}\right)
       \left\{\frac12\left(1-\frac3{2L_0}-\frac1{3L_0^2}\right)
                       -\frac{160000}{X_0^2}\right\}.
                                                    \label{eq:global-A0}
\end{align}
The coefficient on the final line is positive. As $X$ increases, its
braced lower bound increases; the negative logarithmic errors decrease
in modulus. This proves the asserted uniformity in
\eqref{eq:global-arithmetic-growth}. Evaluation gives
\begin{equation}\label{eq:global-A0-value}
 A_0>10.8131594547983.
\end{equation}

\subsection{Uniform bounds for the derivative and radial terms}
Retain only prime powers $q=p^a\leq1000$, $a\geq2$, and put
\[
 \sigma=\sum_{p^a\leq1000,\ a\geq2}\frac{\log p}{p^a},
 \quad\ell_- =\tfrac14\log X_0+\log c_- -C_+-\sigma,
 \quad\ell_+ =\tfrac14\log X_0+\log c_R-C_- -\sigma.
\]
By \eqref{eq:global-geometry}, the actual slope is between
$\ell_-+\frac14\log(X/X_0)$ and
$\ell_++\frac14\log(X/X_0)$. Discarding further factorial savings
gives a valid upper bound for the determinant.

Each residue modulo $q$ can occur in at most $n\lceil J/q\rceil$
selected columns. Dividing by $M$ bounds this capacity by
$a_q=(1+q/J_0)/(\theta_-q)$.
Put $b_q=\min(q,\lfloor1/a_q\rfloor)$ and
$r_q=(1-b_qa_q)_+$. Filling the largest remainders first proves that
\begin{equation}\label{eq:global-column-remainder}
 c_j=2\sum_{q=p^a\leq1000,\ a\geq2}\frac{\log p}{q}
 \left\{a_q\left(b_q(q-1)-\frac{b_q(b_q-1)}2\right)
                       +r_q(q-1-b_q)_+\right\}
       <60.016864219225
\end{equation}
is an admissible bound for the column remainder.

For the factorial saving, use $3200$ intervals $[u,v]$ of equal length
in $[0,1]$. On each interval take
\[
 a=(1-\theta_+(1-u)^2)e_-,\quad
 b=\min(1,(1+1/T)v+1/R_0),\quad
 \alpha=[(1-u)^2-(1-v)^2-1/M_0]_+.
\]
The number of quantiles in this interval is at least $\alpha M$.
The monotonicities following \eqref{eq:radial-rearrangement}, with
$\ell=\ell_-$, therefore give
\begin{equation}\label{eq:global-factorial-saving}
 S_{\rm der}=3\sum_{[u,v]}\alpha h(b,a)
                       >0.1803971924811329636.
\end{equation}
All the inequalities used here improve with $R,S,J$ at their stated
lower bounds; no new derivative orders or prime powers are assumed
to have the same residues as those at the endpoint.

We next put the convex estimate on a fixed domain. Set $N=800$,
$b_0=480$, and $n_0=1\,842\,400$, which is a lower bound for $n$ by
\eqref{eq:global-node-bounds}. In \eqref{eq:radial-bin-weights} use
$\eta_+$ and $n_0$, and denote the resulting numbers by $L_b^0$.
Take
\[
 \overline L_0=\tfrac12\log(N/(b_0+1)),\qquad
 u_0=(J_*+2)/R_0,
\]
and let $\Delta_0$ be \eqref{eq:radial-Delta} evaluated at $(R_0,T,J_*)$.
For every $D\geq D_0$ the actual threshold $(J+2)/R$ is at most
$u_0$. Also $\Delta\leq\Delta_0$: substitute
$(J+1)/R\leq e+1/R_0$ in its quadratic expression, whose derivative
is positive above $1/(2K)$, and compare with $J_*/R_0$.
The resulting strict bound absorbs the omitted negative term $-1/R^2$.

Choose $120$ integers
\[
 m_r=\lfloor.02J_0\rfloor+
 \left\lfloor\frac{r(\lfloor.46J_0\rfloor-\lfloor.02J_0\rfloor)}{119}
                                                      \right\rfloor,
 \quad 0\leq r<120,
\]
and put $\tau_r=m_r/J_0$. Retain the initial segment on which
\[
 G_r=1-\frac{\tau_r+1/J_0}{\theta_-}>0,
 \quad a_r=1-\sqrt{1-G_r}-(1-\tau_r)J_*/R_0>0,
 \quad\nu_r=\ell_-+\log a_r>0.
\]
Use successive differences of these decreasing $\nu_r$, with the next
value after the last retained cut equal to zero, and reduce the cuts
to $\widetilde\tau_r=(\theta_-/\theta_+)\tau_r$.
These choices give lower staircase levels for the cuts
$\lceil\tau_rJ(D)\rceil$ at every $D\geq D_0$; the extra $1/J_0$
accounts for rounding the cuts upward.

In \eqref{eq:radial-functional} replace $\theta$ by $\theta_-$,
the radial coefficients by $L_b^0$, the threshold in $\Phi$ by $u_0$,
and the quadratic coefficient $3\varepsilon/(2\theta)$ by
$3(J_0-1)/(2R_0\theta_-)$. Use the cuts
$\widetilde\tau_r$ and the preceding staircase levels, and replace
the negative correction by $-\overline L_0\Delta_0$.
Call the resulting function $\mathcal F_0$.
For any actual occupancy vector $w$, put
$\widetilde w=(\theta_-/\theta)w$. Then
$\widetilde w\in\mathcal P_{\theta_-}$ and its cumulative fractions
$G_b$ are unchanged. Moreover,
\[
 \frac{\varepsilon}{\theta}w_b^2
 =\frac{\varepsilon\theta}{\theta_-^2}\widetilde w_b^2
 \geq\frac{J_0-1}{R_0\theta_-}\widetilde w_b^2.
\]
The same calculation applies to the positive-part squares, since their
rescaled cuts are at least $\widetilde\tau_r$.
The radial weights and their successive differences decrease when either
$\eta$ or the shift $1/n$ increases; this follows by differentiating
the logarithm in \eqref{eq:radial-bin-weights}.
Finally $\overline L_0$ is an upper bound for the clipping weight.
Each of these replacements therefore weakens the lower bound, and
\begin{equation}\label{eq:global-functional-comparison}
 \mathcal F_D(w)\geq\mathcal F_0(\widetilde w).
\end{equation}

The supporting-plane formula \eqref{eq:radial-supporting-plane}, applied
to $\mathcal F_0$, gives
\begin{equation}\label{eq:global-convex-value}
 \min_{\mathcal P_{\theta_-}}\mathcal F_0
                          >0.0825081841978627281=:F_*.
\end{equation}
The supporting-plane evaluation uses only the explicit function above and
the finite linear minimization in~\eqref{eq:radial-supporting-plane}; the
repository code records the parameter values and evaluates it with
outward-rounded arithmetic.

\subsection{The resulting contradiction}
Put
\[
 A=1+\frac1K+\frac{3(K-1)}{2KT}+\frac{K-5}{2KT^2},\qquad
 B=\frac{3(T+1)J_*}{R_0T}(1-\theta_-/2)
                     +\frac{3\theta_+J_*}{R_0^2T}.
\]
These bound $(W_1+W_2)/W$ and $\gamma$ in
\eqref{eq:radial-derivative-coupled} from above.
Combining the radial determinant bound with the arithmetic column
remainder gives the Archimedean upper bound
\begin{equation}\label{eq:global-size-growth}
 U_0+\tfrac34(A+B)\log(X/X_0),
\end{equation}
where
\begin{align}
 U_0={}&A(\tfrac34\log X_0+C_+)+B\ell_+-F_*-S_{\rm der}\notag\\
 &+\frac{3(T+1)c_j}{2R_0T}
                         +\frac{3(T+1)\log M_0}{R_0T}
       <10.8128420196003.                  \label{eq:global-U0}
\end{align}
The coefficient of $c_j$ decreases with $S$. So does the last term:
$((S+1)/S^2)\log(KS(S+1)/2)$ has negative derivative for $S\geq T$,
as follows by logarithmic differentiation and $\log(KS(S+1)/2)>2$.
The geometry grows by at most $3A\log(X/X_0)/4$, while the derivative
slope grows by at most $B\log(X/X_0)/4$.
The larger coefficient in \eqref{eq:global-size-growth} is thus valid.

Subtracting \eqref{eq:global-size-growth} from
\eqref{eq:global-arithmetic-growth} yields
\begin{equation}\label{eq:tail-positive}
 G_{\rm tail}(D)>0.00031743+0.127829\log(X/X_0)>0.
\end{equation}
Here $A_0-U_0>0.0003174351980$ and
$1-3(A+B)/4>0.1278297612987827$.
The nonzero minor supplied above must satisfy the opposite determinant
inequality. This excludes every $D\geq D_0$.

\subsection{The finite interval below \texorpdfstring{$D_0$}{D0}}
For $875\leq D\leq D_0$, keep $P,K,S,R,J,n$ fixed on each closed
interval $[d_-,d_+]$. The lattice radius is enclosed using its endpoint
values. Choose $n$ to be a multiple of $800$ satisfying the resultant
conditions with the upper radius and the lower value of $D$.
The arithmetic estimates use norm representations when $P>(d_--3)/4$;
otherwise every retained prime is inert.

In all cases use $N=800$ radial groups, $3200$ groups for the factorial
saving, and $120$ cuts in \eqref{eq:radial-derivative-coupled}.
The latter run from $\max(1,\lfloor.02J\rfloor)$ to
\[
 \max\left(\max(1,\lfloor.02J\rfloor)+1,
       \min(\lfloor M/n\rfloor-1,\lfloor.495J\rfloor)\right),
\]
equally spaced with integer rounding, retaining the initial segment
where the stated positivity conditions hold. Take $b_0=480$.
The prime-power cutoff is $J-1$.

The intervals using norm representations have the parameter families
in Table~\ref{tab:global-hybrid-parameters}; there
$S=\lfloor c_RP/K\rfloor$, $R=KS$, $J=\lfloor eR\rfloor$.
Their slope windows have cutoff at most $750000$ in the first row and
$140000$ thereafter. The prime-product bounds in Appendix~\ref{app:arithmetic}
apply on each constituent interval.
\begin{table}[ht]
\centering\small
\begin{tabular}{lrrrr}
\toprule
Range & $P$ & $K$ & $c_R$ & $e$\\
\midrule
$[875,1167]$ & $8\,999\,999$ & $23$ & $7/5$ & $7/125$\\
$[1167,1359]$ & $1\,699\,999$ & $23$ & $8/5$ & $27/500$\\
$[1359,1871]$ & $1\,699\,999$ & $23$ & $9/5$ & $27/500$\\
$[1871,6959]$ & $1\,699\,999$ & $25$ & $2$ & $7/125$\\
$[6959,319591]$ & $1\,699\,999$ & $29$ & $11/5$ & $29/500$\\
$[319591,3968539]$ & $1\,699\,999$ & $31$ & $23/10$ & $29/500$\\
\bottomrule
\end{tabular}
\caption{Parameter families for the norm-form comparisons. Each range is
subdivided into closed intervals on which all integer parameters are fixed.}
\label{tab:global-hybrid-parameters}
\end{table}

For the remaining intervals take $K=31$, $c_R=23/10$, $e=29/500$, and
$P=\lfloor(d_--3)/4\rfloor$. Between $3968539$ and $4894699$ the prime
sums are evaluated directly. Above $4894699$ use the explicit bounds of
\cite[Theorems 4, 6 and 18]{rosserschoenfeld}; the additional estimate
$\vartheta(x)<x$ is used only for $x\leq10^8$, within its stated range.
On each interval the radial function is bounded below with the upper
value of $\eta$, while the negative clipping term uses the lower value.
The derivative scale likewise uses the appropriate endpoint in each
term. These substitutions bound the entire interval.

The adjoining intervals cover $[875,D_0]$, and their lower bounds are
summarized in Table~\ref{tab:coverage}. Each is positive. Together with
\eqref{eq:tail-positive}, this excludes every $D\geq875$.
The interval endpoints and the numerical calculations are available with the
\href{https://github.com/dualverse-ai/tenth_discriminant/tree/main/bound-498/code}{accompanying code}.
